\section{Introduction}\label{sec:intro}

Byzantine fault-tolerant (BFT) consensus protocols underpin modern
distributed systems, from replicated state machines~\cite{CastroLiskov1999}
to blockchain platforms~\cite{GiladHMVZ2017, YinMGRGA2019}.  Getting these
protocols right is extraordinarily difficult.  Subtle bugs in threshold
conditions, message-handling logic, or view-change sub-protocols have
caused real-world safety violations and costly exploits.  Formal
verification can catch such bugs before deployment---but existing tools
force protocol designers into an uncomfortable trade-off between
\emph{expressiveness} and \emph{scalability}.

On one end of the spectrum, \emph{explicit-state model checkers} like
SPIN~\cite{Holzmann2003} and TLC (for TLA+~\cite{LamportTLA2002})
enumerate concrete system configurations.  Their results are inherently
bounded: verifying a protocol with $n{=}4$ processes says nothing about
$n{=}100$.  On the other end, \emph{parameterized} techniques based on
threshold automata and counter
abstraction~\cite{KonnovVW2015, KonnovLVW2017} treat the number of
processes~$n$, the fault tolerance~$t$, and the actual number of
faults~$f$ as symbolic parameters.  A single safety proof then covers
\emph{all} system sizes satisfying a resilience condition such as
$n > 3t$.  The ByMC tool~\cite{KonnovWidder2018} demonstrated the
viability of this approach on simplified BFT protocols.

However, three practical challenges remain unaddressed by existing
threshold-automata tools:

\begin{enumerate}
\item \textbf{Abstraction fidelity.}
  Classical counter abstraction collapses all processes in the same
  protocol location into a single aggregate counter.  While sound for
  safety (a SAFE result transfers to the real protocol), this
  over-approximation loses per-process message-delivery semantics:
  sender identities, distinct-sender thresholds (``receive $\geq 2t{+}1$
  \emph{distinct} Prepare messages''), and equivocation control become
  invisible.  A SAFE verdict may therefore mask real bugs that only
  manifest under faithful message delivery.

\item \textbf{Cryptographic objects.}
  Modern BFT protocols rely on quorum certificates, threshold
  signatures, and signed messages whose non-forgeability guarantees are
  integral to safety arguments.  No existing threshold-automata tool
  encodes these as first-class semantic objects.

\item \textbf{Trust in the verifier.}
  An SMT-based model checker is a complex artifact.  A bug in the
  parser, the lowering pass, or the SMT encoding could produce a
  spurious SAFE verdict.  Without an independent mechanism to validate
  the proof, the user must trust the entire toolchain.
\end{enumerate}

\paragraph{Contributions.}
We present \tarsier{}, a parameterized model checker for threshold-guarded
distributed protocols that addresses all three challenges.  Our
contributions are:

\begin{itemize}
\item A \textbf{hierarchy of four network abstraction modes}
  (Sect.~\ref{sec:semantics})---\emph{classic}, \emph{identity-selective},
  \emph{cohort-selective}, and \emph{process-selective}---with formally
  characterized soundness relationships.  The most precise mode,
  process-selective, is \emph{instance-exact}: its reachable states
  coincide with those of the faithful protocol model, so both SAFE
  \emph{and} UNSAFE verdicts transfer.

\item \textbf{First-class cryptographic objects}
  (Sect.~\ref{sec:language})---certificates, threshold signatures, and
  signed-message channels---with non-forgeability, distinct-signer
  thresholds, and conflict constraints encoded directly as SMT
  assertions.

\item A \textbf{minimal proof kernel} (Sect.~\ref{sec:verification})
  with only four library dependencies that validates proof certificates
  independently of the verification engine, following the
  proof-carrying-code paradigm~\cite{Necula1997}.

\item A \textbf{domain-specific language} (Sect.~\ref{sec:language})
  designed to be close to the pseudo-code in protocol papers while
  remaining formal enough for automated verification, together with
  a deterministic lowering pipeline to threshold automata.

\item An \textbf{evaluation} (Sect.~\ref{sec:evaluation}) on
  49~protocol models from five major BFT families (PBFT, HotStuff,
  Tendermint, Paxos, Algorand) and nine additional protocols.  We show
  that faithful network modes close real soundness gaps while remaining
  tractable, and that proof certificates reduce the trusted computing
  base to four library dependencies.
\end{itemize}

\tarsier{} is implemented in Rust as a modular workspace of 11~crates,
using Z3~\cite{DeMoura2008} and cvc5~\cite{BarbosaBBKLMMMN2022} as
SMT backends.
